%% my-new-words-env.tex --------------------------------------------
%% 定义环境 MyNewWords：用于编写“简单的”生词本。
%%
%% 单词条目的编写格式（三种任选）：
%%   \i 单词\w 音标\s 释义\e
%%   \i 单词\w 释义\e
%%   \i 单词\w\e
%%
%% 环境的可选参数由 pgfkeys 处理，目前支持：
%%   sort=true  按单词的 ASCII 编码对条目排序（大写在前）
%%   sort=false 保持书写顺序（默认）
%%
%% 音标按 tipa 的“简写”方式书写，环境会自动把它放进 \textipa，
%% 例如：\i word\w ""w3:d\s ... \e
%%
%% 本文件需要以下宏包（若已在导言区加载 arraysort，
%% 请务必带上 comparestr 选项：\usepackage[comparestr]{arraysort}）：
%%   pgfkeys、tipa、arraysort[comparestr]
%%
%% 注意：
%%   * 每个条目必须以 \e 结束；
%%   * \i、\w、\s、\e 是环境的专用分隔符，不要在正文中另作它用
%%     （\s 只能出现在音标与释义之间一次）。
\ProvidesFile{my-new-words-env.tex}%
  [2026/08/24 v1.0 简单的生词本环境 MyNewWords]

\makeatletter

% ---------- 依赖宏包 ----------
\IfPackageLoadedF{pgfkeys}{\RequirePackage{pgfkeys}}
\IfPackageLoadedF{tipa}{\RequirePackage{tipa}}
\IfPackageLoadedTF{arraysort}
  {\@ifpackagewith{arraysort}{comparestr}{}{%
     \PackageWarningNoLine{my-new-words}{%
       arraysort 已加载但缺少 comparestr 选项,\MessageBreak
       按字符串排序可能不正确}}}
  {\PassOptionsToPackage{comparestr}{arraysort}%
   \RequirePackage{arraysort}}

% ---------- 状态 ----------
\newif\ifMNW@sort
\newcount\MNW@n              % 已收集的条目数
\newcommand*\MNW@index{}     % 排序数组元素中的索引标记
\newcommand*\MNW@stop{}      % 参数扫描哨兵

% 建立排序用数组（仅一次），每个环境开始时把长度清零
\newarray\MNW@arr
\def\MNW@initarr{\expandafter\gdef\csname total@MNW@arr\endcsname{0}}

% ---------- 选项 ----------
\pgfkeys{
  /MyNewWords/.cd,
  sort/.is if = MNW@sort,
  sort/.default = true,
  unknown/.code = {%
    \PackageWarning{my-new-words}{忽略未知选项 `\pgfkeyscurrentkey'}}%
}

% ---------- 条目解析 ----------
% \i <单词>\w <其余>\e
\long\def\MNW@entry#1\w#2\e{%
  \global\advance\MNW@n\@ne
  \edef\MNW@cur{\the\MNW@n}%
  \edef\MNW@tmp{MNW@word@\MNW@cur}%
  \expandafter\gdef\csname\MNW@tmp\endcsname{#1}%
  % 追加三个哨兵 \s：其中两个充当 #1、#2 的定界符，
  % 剩下的连同数据中多余的 \s 一并落入弃置参数 #3，
  % 因此无论条目里有没有 \s 都能正确切分
  \MNW@split#2\s\s\s\MNW@stop
  % 存入排序数组：元素 = <单词>\MNW@index<序号>
  \expandafter\MNW@store\expandafter{\csname MNW@word@\MNW@cur\endcsname}}

\long\def\MNW@store#1{\MNW@arr(\MNW@cur)={#1\MNW@index\MNW@cur}}

% <其余> 切分为 音标 与 释义（可缺省）；#3 吞掉多余的哨兵
\long\def\MNW@split#1\s#2\s#3\MNW@stop{%
  \def\MNW@pb{#2}%
  \ifx\MNW@pb\@empty
    \MNW@saveparts{}{#1}%        无 \s：#1 是释义，没有音标
  \else
    \MNW@saveparts{#1}{#2}%
  \fi}

\long\def\MNW@saveparts#1#2{%
  \edef\MNW@tmp{MNW@phon@\MNW@cur}%
  \expandafter\gdef\csname\MNW@tmp\endcsname{#1}%
  \edef\MNW@tmp{MNW@mng@\MNW@cur}%
  \expandafter\gdef\csname\MNW@tmp\endcsname{#2}}

% ---------- 输出 ----------
\long\def\MNW@emitone#1{%
  \item {\bfseries\csname MNW@word@#1\endcsname}%
  \expandafter\ifx\csname MNW@phon@#1\endcsname\@empty\else
    \quad\expandafter\textipa\expandafter{\csname MNW@phon@#1\endcsname}%
  \fi
  \par
  \expandafter\ifx\csname MNW@mng@#1\endcsname\@empty\else
    \csname MNW@mng@#1\endcsname\par
  \fi}

\def\MNW@loop@plain{%
  \count@=\z@
  \@whilenum\count@<\MNW@n\do{%
    \advance\count@\@ne
    \MNW@emitone{\the\count@}}}

\def\MNW@loop@sorted{%
  \count@=\z@
  \@whilenum\count@<\MNW@n\do{%
    \advance\count@\@ne
    \testarray{MNW@arr}(\the\count@)%
    \expandafter\MNW@pickindex\temp@macro\MNW@stop}}

\long\def\MNW@pickindex#1\MNW@index#2\MNW@stop{\MNW@emitone{#2}}

\def\MNW@output{%
  \ifnum\MNW@n>0
    \begin{list}{}{%
      \setlength{\leftmargin}{0pt}%
      \setlength{\itemindent}{0pt}%
      \setlength{\listparindent}{0pt}%
      \setlength{\itemsep}{0.8em}%
      \setlength{\parsep}{0.4em}%
      \setlength{\topsep}{0.5em}}%
      \ifMNW@sort
        \ifnum\MNW@n>\@ne \sortArray{\the\MNW@n}{MNW@arr}\fi
        \MNW@loop@sorted
      \else
        \MNW@loop@plain
      \fi
    \end{list}%
  \fi}

% ---------- 环境 ----------
\newenvironment{MyNewWords}[1][]{%
  \MNW@sortfalse
  \def\MNW@opts{#1}%
  \ifx\MNW@opts\@empty\else\pgfkeys{/MyNewWords/.cd,#1}\fi
  \global\MNW@n=\z@
  \MNW@initarr
  \parskip=\z@
  \let\i\MNW@entry              % 条目起始；\w、\s、\e 仅作扫描定界符
  \ignorespaces
}{%
  \MNW@output
}

\makeatother
\endinput
